\section{Geometric c/o structures}
\label{geosection}

In \cite{KP} we axiomatized a structure of spaces that allow four dif\/ferent types of gluing: i.e.\
either open or closed and either self or non-self gluings together with a coloring of the open gluings by brane labels~$\beta$. In the above case the gluing is on two  windows, which are either both
open or closed and either belong to the same (self gluing) or disjoint surfaces (non-self gluing).
The coloring is by the brane labels. The color of a window is the pair of brane labels of its
end points. When we glue $\beta$ arc families, we glue the surfaces and
glue the arc families as foliations. We will now review this process according to~\cite{KP}.


\subsection{The gluing underlying the topological c/o structure}




If $\alpha\in\widetilde{\rm Arc} (n,m)$, then def\/ine the {\it $\alpha$-weight} $\alpha(w)$ of an active window $w$ to be the sum
of
the weights of arcs in $\alpha$ with endpoints in $w$, where we count with multiplicity
(so if an arc in $\alpha$ has both endpoints in $w$, then the weight of this arc contributes twice to the weight of
$w$).

Suppose we have a pair of arc families $\alpha_1$, $\alpha _2$
in respective windowed surfaces $F_1$, $F_2$ and a~pair of active windows $w_1$ in $F_1$ and $w_2$ in $F_2$, so that the
$\alpha _1$-weight of $w_1$ agrees with the $\alpha _2$-weight of $w_2$.
Since $F_1$, $F_2$ are oriented surfaces, so are the windows $w_1$, $w_2$ oriented.   In each operation, we
identify windows reversing orientation, and we identify certain distinguished points.

To def\/ine the open and closed gluing ($F_1\neq F_2$) and self-gluing ($F_1=F_2$) of $\alpha_1$, $\alpha _2$ along the
windows
$w_1$, $w_2$, we
identify windows and distinguished points in the natural way and combine foliations.

A crucial dif\/ference  between the closed and open string operations is that in the closed case,
the points are thought of as marked, which in the open case the points behave like punctures.
This means that in the closed case, we
replace the  distinguished point by simply forgetting that is was distinguished. This way  no puncture is created. In the open case the distinguished
points
always give rise to other distinguished points or perhaps punctures.
 In any case whenever distinguished points are
identif\/ied, one takes the union of brane labels (the intersection of branes) at the new resulting distinguished point or puncture.

\begin{figure}[th]
%\epsfxsize = \textwidth \epsfbox{stropfig4.eps}
\centerline{\includegraphics{Kaufmann-Fig1}}

\caption{The dif\/ferent cases of gluing.}\label{stropfig4}
\end{figure}


\subsubsection{Closed gluing and self-gluing}
Identify the two corresponding boundary components of $F_1$ and
$F_2$, identifying also the distinguished points on them {\it and
then including this point in the resulting surface} $F_3$. $F_3$~inherits a~brane-labeling from those on $F_1$, $F_2$ in the natural
way.  We furthermore glue $\alpha _1$ and $\alpha _2$ together in
the natural way, where the two collections of foliated rectangles
in~$F_1$ and~$F_2$ which meet $w_1$ and $w_2$ have the same total
width by hypothesis and therefore glue together naturally to
provide a measured foliation ${\mathcal F}$ of a closed subsurface
of $F_3$.


\subsubsection{Open gluing}

The surfaces $F_1$ and $F_2$ are distinct, and we identify $w_1$ to $w_2$ to produce
$F_3$.
There are cases depending upon whether the closure of $w_1$ and $w_2$ is an interval or a circle.  The salient
cases are
illustrated in Fig.~\ref{stropfig4},~b--d.  In each case,  distinguished points on the boundary in~$F_1$ and~$F_2$ are
identif\/ied to produce a new distinguished boundary point in $F_3$, and the brane labels are combined, as is also
illustrated.  As before, since the
$\alpha _1$-weight on $w_1$ agrees with the $\alpha _2$-weight on $w_2$, the foliated rectangles again combine to
provide a measured foliation~${\mathcal F}$ of a closed subsurface of~$F_3$.

 {\bf Open self-gluing.}~There are again cases depending upon whether the closure of $w_1$ or $w_2$ is a~circle or an interval, but there is a further case as well when the two intervals lie in a common boundary component
and
are consecutive.  Other than this last case, the construction is identical to those illustrated
in Fig.~\ref{stropfig4},~b--d.  In case the two windows are consecutive along a common boundary component, again they are identif\/ied
so
as to produce a surface~$F_3$ with a puncture resulting from their common endpoint as in Fig.~\ref{stropfig4},~e--f, where the
puncture is
brane-labeled by the label of this point, and the foliated rectangles combine to provide a measured foliation~${\mathcal F}$ of a closed subsurface of~$F_3$.



At this stage, we have only constructed a measured foliation ${\mathcal F}$ of a closed subsurface of $F_3$, and
indeed,
${\mathcal F}$ will typically not be a weighted arc family, but the sub-foliation
${\mathcal
F}'$ comprised of leaves that meet $\partial F$ corresponds to a weighted arc family $\alpha _3$ in $F_3$.
Notice that
the $\alpha _3$-weight of any window uninvolved in the operation agrees with its $\alpha _1$- or $\alpha _2$-weight.

The assignment of $\alpha_3$ in $F_3$ to $\alpha _i$ in $F_i$, for
$i=1,2$ completes the def\/inition of the various operations.
Associativity and equivariance for bijections are immediate, and so
we have our f\/irst non-trivial example of a c/o structure; see
Appendix~\ref{appendixA} for the precise def\/inition.

